% compile with XeLaTeX
\documentclass[dvipsnames,mathserif]{beamer}
\setbeamertemplate{footline}[frame number]
\setbeamercolor{footline}{fg=black}
\setbeamerfont{footline}{series=\bfseries}
\usepackage{tikz}
\usepackage{xcolor}
\usepackage{booktabs}

\usetheme{Darmstadt}

\makeatletter
\newcommand{\RTListe}{\raggedleft\rightskip\leftm}
\newcommand{\leftm}{\@totalleftmargin}
\makeatother

\setbeamertemplate{frametitle}
{\vspace*{-1mm}
\nointerlineskip
    \begin{beamercolorbox}[sep=0.3cm,ht=2.2em,wd=\paperwidth]{frametitle}
        \vbox{}\vskip-2ex%
        \strut\hskip1ex\insertframetitle\strut
        \vskip-0.8ex%
    \end{beamercolorbox}
}
\makeatletter
\setbeamertemplate{subsection in toc}
{\leavevmode\rightskip=5ex%
\llap{\raise0.1ex\beamer@usesphere{subsection number projected}{bigsphere}\kern1ex}%
\inserttocsubsection\par%
}
\makeatother

\setbeamertemplate{itemize item}{\scriptsize\raise1.25pt\hbox{\donotcoloroutermaths$\blacktriangleleft$}} 

\defbeamertemplate{enumerate item}{square2}
{\LR{
    \hbox{%
    \usebeamerfont*{item projected}%
    \usebeamercolor[bg]{item projected}%
    \vrule width2.25ex height1.85ex depth.4ex%
    \hskip-2.25ex%
    \hbox to2.25ex{%
    \hfil%
    {\color{fg}\insertenumlabel}%
    \hfil}%
}%
}}

\setbeamertemplate{enumerate item}[square2]
\setbeamertemplate{navigation symbols}{}

% FOSSEE LOGO ON EVERY PAGE
\logo{\includegraphics[height=0.9cm]{fossee_logo.png}}

\setbeamertemplate{caption}[numbered]
\begin{document}

\rightskip\rightmargin
\title{\textbf{Unstructured Mesh Generation} \\ Technical Progress Report (Week 11)}
\author{\Large \textbf{Pranjal Kumar Shukla}}
\institute{\large\textbf{FOSSEE Internship -- Progress Report}\\[4pt]
SCAI School\\[2pt]
VIT Bhopal University}
\footnotesize{\date{\today}}


\begin{frame}
\maketitle
\end{frame}

%% ─────────────────────────────────────────
%% 1. WEEK 11 FOCUS
%% ─────────────────────────────────────────
\section{Week 11 Focus}
\begin{frame}{Week 11 Focus: Stability, UX \& Correctness}
    \begin{alertblock}{Objective}
        No new features added this week. Full focus on \textbf{improving robustness, fixing bugs, and enhancing the user experience} of existing tools.
    \end{alertblock}
    \vspace{0.3cm}
    \textbf{Key improvement areas:}
    \begin{itemize}
        \item Geometric validation and smart error messages
        \item Performance optimization (no UI freezes)
        \item Metric unit display on all length inputs
        \item Stable multi-frame CFD animation pipeline
    \end{itemize}
\end{frame}

%% ─────────────────────────────────────────
%% 2. GEOMETRIC VALIDATION
%% ─────────────────────────────────────────
\section{Improvements}
\begin{frame}{1. Geometric Validation for Trailing Edge}
    \textbf{Problem:} Trailing Edge Refinement was silently applied to all objects, including spheres -- where it is physically meaningless.
    \vspace{0.2cm}
    \begin{block}{Solution: \texttt{bmesh} Edge Scan}
        \begin{itemize}
            \item Scans all edges using the \texttt{bmesh} API before meshing
            \item Computes face-face angle for each edge
            \item If \textbf{no edge exceeds 30°}: shows a clear error message
            \item Allows objects with \textbf{both curved and sharp edges} (e.g.\ cylinders with rims)
        \end{itemize}
    \end{block}
    \vspace{0.1cm}
    \textit{``Trailing edge refinement requires sharp edges. Your geometry appears completely smooth (e.g.\ a sphere). Please disable this feature.''}
\end{frame}

\begin{frame}{2. Non-Manifold Geometry Pre-Check}
    \textbf{Problem:} cfMesh silently fails or produces corrupt meshes when the input STL has holes or self-intersections.
    \vspace{0.3cm}
    \begin{block}{Solution: Pre-Meshing Validation}
        \begin{itemize}
            \item \texttt{bmesh} scan detects non-manifold edges before export
            \item Reports exact edge count: \textit{``X non-manifold edges detected -- please fix in Edit Mode''}
            \item Non-blocking warning (cfMesh may still succeed on simple cases)
        \end{itemize}
    \end{block}
\end{frame}

\begin{frame}{3. Performance: Triangulation Optimization}
    \textbf{Problem:} Previous implementation used \texttt{bpy.ops.object.modifier\_add()} in a loop, triggering a full Blender dependency graph rebuild per object. This caused \textbf{UI freezes} on complex meshes.
    \vspace{0.3cm}
    \begin{block}{Solution: Low-Level \texttt{bmesh} Operations}
        \begin{itemize}
            \item Replaced with \texttt{bmesh.ops.triangulate(bm, faces=bm.faces[:])}
            \item Operates directly on mesh data -- no graph rebuilds
            \item Significantly faster on dense meshes
        \end{itemize}
    \end{block}
    \vspace{0.1cm}
    \begin{center}
        \small
        \begin{tabular}{lcc}
            \toprule
            \textbf{Method} & \textbf{Graph Rebuild} & \textbf{Speed} \\
            \midrule
            \texttt{bpy.ops} loop & Per object & Slow \\
            \texttt{bmesh.ops} & None & Fast \\
            \bottomrule
        \end{tabular}
    \end{center}
\end{frame}

\begin{frame}{4. Beginner-Friendly Error Messages}
    \textbf{Problem:} Errors surfaced raw Python exceptions, which are meaningless to non-developers.
    \vspace{0.2cm}
    \begin{block}{Improved Exception Handling (3 files)}
        \begin{itemize}
            \item \textbf{ops\_meshing.py}: \texttt{PermissionError} $\rightarrow$ ``Check write permissions''; \texttt{MemoryError} $\rightarrow$ ``Increase cell size''
            \item \textbf{ops\_geometry.py}: STL import errors now guide users to valid ASCII/Binary STL files
            \item \textbf{ops\_inspect.py}: VTU parse error now says ``Did you run the mesher first?''
            \item \textbf{OpenFOAM guard}: Checks if \texttt{cartesianMesh} is accessible before starting any operation
        \end{itemize}
    \end{block}
\end{frame}

%% ─────────────────────────────────────────
%% 3. UI IMPROVEMENTS
%% ─────────────────────────────────────────
\section{UI Improvements}
\begin{frame}{5. Metric Units \& Live Estimators in UI}
    \begin{columns}
        \column{0.5\textwidth}
        \textbf{Metric Units}
        \begin{itemize}
            \item Added \texttt{subtype=`DISTANCE'} to all length-based properties
            \item Blender now auto-displays ``m'', ``cm'', ``mm'' next to inputs
            \item Applies to: cell size, layer thickness, characteristic length, trailing edge size
        \end{itemize}
        \column{0.5\textwidth}
        \textbf{Live Cell Count Estimator}
        \begin{itemize}
            \item Before meshing: displays \textit{``Est. Cells: $\sim$5,000''}
            \item Formula: $n \approx V_{\text{bbox}} / dx^3$
            \item Warns in red if $> 2\times10^6$ cells
            \item Helps users avoid out-of-memory crashes
        \end{itemize}
    \end{columns}
\end{frame}

\begin{frame}{6. Courant Number Warning System}
    \textbf{Problem:} icoFoam diverges when the Courant number $\text{Co} = U \cdot \Delta t / \Delta x > 1$. Users had no feedback before running.
    \vspace{0.2cm}
    \begin{block}{Live Courant Checker (Physics \& Solver Panel)}
        \[
            \text{Co} = \frac{U_{\text{mag}} \cdot \Delta t}{\Delta x}, \quad \Delta t_{\text{safe}} = \frac{0.5 \cdot \Delta x}{U_{\text{mag}}}
        \]
        \begin{itemize}
            \item Updates live as you change velocity, $\Delta t$, or cell size
            \item \textcolor{green}{\textbf{Green}}: Co $< 0.5$ (stable)
            \item \textcolor{orange}{\textbf{Orange}}: $0.5 <$ Co $< 1.0$ (borderline)
            \item \textcolor{red}{\textbf{Red alert}}: Co $\geq 1$ (solver WILL diverge) + shows safe $\Delta t$
        \end{itemize}
    \end{block}
\end{frame}

\begin{frame}{7. Estimated Animation Frames in Time Controls}
    \textbf{Problem:} Users did not know how many output files would be written before running a simulation.
    \vspace{0.2cm}
    \begin{block}{Live Frame Count Preview}
        \textbf{icoFoam} (adjustableRunTime):
        \[
            n_{\text{frames}} = \texttt{write\_interval} + 1
        \]
        \textbf{simpleFoam} (step-based):
        \[
            n_{\text{frames}} = \lfloor \texttt{max\_iterations} / \texttt{write\_interval} \rfloor
        \]
        \begin{itemize}
            \item Shown live inside the Time \& Interval box
            \item Warns if 0 or 1 frame (too coarse interval)
            \item Warns if $>500$ frames (disk space risk)
        \end{itemize}
    \end{block}
\end{frame}

%% ─────────────────────────────────────────
%% 4. ANIMATION PIPELINE FIXES
%% ─────────────────────────────────────────
\section{Animation \& Post-Processing}
\begin{frame}{8. Animation Pipeline: Multi-Frame Fix}
    \textbf{Problem:} Animation showed only 1 visible frame regardless of time steps.
    \vspace{0.2cm}
    \begin{block}{Root Cause}
        Old keyframe logic set each mesh visible for exactly 1 Blender frame, then immediately hidden again:
        \textit{frame N-1: hidden, frame N: visible, frame N+1: hidden}
    \end{block}
    \begin{block}{Fix}
        \begin{itemize}
            \item Each mesh visible from its frame until the \textbf{next one appears}
            \item \texttt{CONSTANT} keyframe interpolation (no blending between steps)
            \item Compatible with Blender 3.x / 4.x / 5.x action API
            \item Last frame stays visible permanently
        \end{itemize}
    \end{block}
\end{frame}

\begin{frame}{9. Parallel Run Reconstruction (\texttt{reconstructPar})}
    \textbf{Problem:} When using $>1$ CPU cores, OpenFOAM writes results to \texttt{processor0/}, \texttt{processor1/}..., \textbf{not} the main case directory. \texttt{foamToVTK} saw only the initial condition, producing just 1--2 frames.
    \vspace{0.2cm}
    \begin{block}{Automatic Fix in ``Color Boundary Mesh''}
        \begin{enumerate}
            \item Detects \texttt{processor*} directories automatically
            \item Runs \texttt{reconstructPar -newTimes} before \texttt{foamToVTK}
            \item Only merges new time steps (fast on re-run)
            \item Falls back gracefully with a warning if reconstruction fails
        \end{enumerate}
    \end{block}
    \textbf{Result:} 21 frames correctly generated from 20 write intervals.
\end{frame}

\begin{frame}{10. Other Post-Processing Improvements}
    \begin{itemize}
        \item \textbf{Auto Shade Smooth}: Imported result STL is automatically smooth-shaded, eliminating the ``faceted sphere'' appearance without manual steps
        \item \textbf{Mesh Quality Icons}: Non-ortho, skewness, and aspect ratio fields each display {\color{green}$\checkmark$} / {\color{orange}$\bigcirc$} / {\color{red}$\times$} icons based on OpenFOAM thresholds
        \item \textbf{Octree Exact-Multiple Bug}: \texttt{maxCellSize} is nudged by $10^{-5}$ relative epsilon before writing \texttt{meshDict}, preventing the \textit{``0 cells in mesh''} crash
        \item \textbf{Live Log Tail}: Status box updates in real-time with the last cfMesh/OpenFOAM output line during execution
    \end{itemize}
\end{frame}

%% ─────────────────────────────────────────
%% 5. TIMELINE
%% ─────────────────────────────────────────
\section{Timeline}
\begin{frame}{Updated Timeline \& Milestones}
    \begin{table}
    \centering
    \footnotesize
    \begin{tabular}{|c|p{6.5cm}|l|}
    \hline
    \textbf{Week} & \textbf{Task} & \textbf{Status} \\ \hline
    4 & Turbulence models (RANS), Exception handling & \textcolor{green}{\textbf{Done}} \\ \hline
    5 -- 6 & STL Importer, Jet Colormaps, Residuals \& Y+ & \textcolor{green}{\textbf{Done}} \\ \hline
    7 & System file opener, Custom liquid presets \& Smart UI & \textcolor{green}{\textbf{Done}} \\ \hline
    8 & OpenFOAM 2412 migration, Modular Architecture & \textcolor{green}{\textbf{Done}} \\ \hline
    9 -- 10 & Real-time Vis, VTK Slices, icoFoam \& simpleFoam & \textcolor{green}{\textbf{Done}} \\ \hline
    \textbf{11} & \textbf{Stability, UX Polish, Animation, Bug Fixes} & \textcolor{green}{\textbf{Done}} \\ \hline
    12 & Final Documentation \& Project Handover & Upcoming \\ \hline
    \end{tabular}
    \end{table}
\end{frame}

%% ─────────────────────────────────────────
%% 6. SUMMARY
%% ─────────────────────────────────────────
\section{Summary}
\begin{frame}{Week 11 Summary}
    \begin{center}
        \small
        \begin{tabular}{p{4.2cm}p{5.8cm}}
            \toprule
            \textbf{Improvement} & \textbf{Outcome} \\
            \midrule
            Geometric Validation & Trailing edge only on sharp-edged geometry \\
            Non-Manifold Check & Warns before meshing broken STL \\
            Triangulation Speed & No UI freeze on large meshes \\
            Beginner Error Messages & Plain-English guidance instead of Python tracebacks \\
            Metric Units & ``m'' / ``mm'' shown on all distance inputs \\
            Courant Warning & Live Co number; prevents icoFoam divergence \\
            Frame Count Estimator & Accurate preview in Time \& Interval panel \\
            Animation Fix & All time steps visible; CONSTANT keyframes \\
            reconstructPar & Parallel runs now produce full animation \\
            Octree Bug Fix & No more ``0 cells'' crash \\
            \bottomrule
        \end{tabular}
    \end{center}
\end{frame}

\begin{frame}
    \begin{center}
        \Huge \textbf{Thank You!} \\
        \vspace{0.5cm}
        \large End of Week 11 Report
    \end{center}
\end{frame}

\end{document}
